% =====================================================================
%  Workshop 6 -- Testing & CI/CD
%  Web Applications -- Semester II 2026
%
%  ESQUELETO DEL REPORTE
%
%  1. Compila TAL CUAL, sin imagenes. Cada captura es una caja gris con
%     el nombre del archivo que espera y la descripcion de lo que debe
%     mostrar.
%
%  2. Para insertar una captura basta con guardarla en  img/  con el
%     nombre exacto de la caja (ej.  img/fig01-endpoint-code.png ) y
%     recompilar. Extensiones aceptadas: .png, .jpg, .jpeg, .pdf.
%
%  3. Todo lo marcado con \todo{...} (en rojo) es texto pendiente que se
%     redacta al final, cuando se pida.
%
%  4. Compilar con pdfLaTeX (Overleaf: New Project -> Upload Project,
%     subiendo este .tex y la carpeta img/).
% =====================================================================

\documentclass[10pt,a4paper]{article}

\usepackage[utf8]{inputenc}
\usepackage[T1]{fontenc}
\usepackage[english]{babel}
\usepackage[margin=2.1cm]{geometry}
\usepackage{graphicx}
\usepackage{xcolor}
\usepackage{listings}
\usepackage{booktabs}
\usepackage{tabularx}
\usepackage{float}
\usepackage{caption}
\usepackage{enumitem}
\usepackage[hidelinks]{hyperref}

\setlist{nosep,leftmargin=1.1cm}
\setlength{\parskip}{0.35em}

% ---------------------------------------------------------------------
%  Colores y estilos de codigo
% ---------------------------------------------------------------------
\definecolor{codebg}{RGB}{248,248,250}
\definecolor{codeframe}{RGB}{212,212,218}
\definecolor{codekey}{RGB}{0,84,166}
\definecolor{codestr}{RGB}{163,21,21}
\definecolor{codecom}{RGB}{106,115,125}
\definecolor{phfill}{RGB}{246,246,248}
\definecolor{phline}{RGB}{170,170,178}
\definecolor{todored}{RGB}{200,30,30}

\lstdefinestyle{base}{
  backgroundcolor=\color{codebg},
  frame=single,
  rulecolor=\color{codeframe},
  basicstyle=\ttfamily\scriptsize,
  keywordstyle=\color{codekey}\bfseries,
  stringstyle=\color{codestr},
  commentstyle=\color{codecom}\itshape,
  showstringspaces=false,
  breaklines=true,
  columns=fullflexible,
  keepspaces=true,
  xleftmargin=5pt,
  xrightmargin=5pt,
  aboveskip=6pt,
  belowskip=6pt,
}

\lstdefinelanguage{yaml}{
  keywords={true,false,null},
  sensitive=false,
  comment=[l]{\#},
  morestring=[b]',
  morestring=[b]",
}

\lstdefinestyle{py}{style=base,language=Python}
\lstdefinestyle{yml}{style=base,language=yaml}
\lstdefinestyle{sh}{style=base,language=bash}
\lstdefinestyle{plain}{style=base,language=}
\lstset{style=plain}

% ---------------------------------------------------------------------
%  Texto pendiente
% ---------------------------------------------------------------------
\newcommand{\todo}[1]{{\color{todored}\textbf{[TODO:} #1\textbf{]}}}

% ---------------------------------------------------------------------
%  Marcador de figura: si la imagen existe la inserta, si no dibuja
%  una caja vacia con la descripcion de la captura que falta.
% ---------------------------------------------------------------------
\newcommand{\phbox}[4]{%
  \begingroup
  \setlength{\fboxsep}{0pt}%
  \fcolorbox{phline}{phfill}{%
    \begin{minipage}[c][#2][c]{\dimexpr#1-2\fboxrule\relax}
      \centering
      {\ttfamily\scriptsize\color{phline}[ img/#3 ]}\\[0.7em]
      \begin{minipage}{0.88\linewidth}
        \centering\scriptsize\itshape #4
      \end{minipage}
    \end{minipage}%
  }%
  \endgroup
}

% #1 ancho  #2 alto de la caja  #3 nombre base del archivo  #4 descripcion
\newcommand{\shot}[4]{%
  \IfFileExists{img/#3.png}{\includegraphics[width=#1]{img/#3.png}}{%
  \IfFileExists{img/#3.jpg}{\includegraphics[width=#1]{img/#3.jpg}}{%
  \IfFileExists{img/#3.jpeg}{\includegraphics[width=#1]{img/#3.jpeg}}{%
  \IfFileExists{img/#3.pdf}{\includegraphics[width=#1]{img/#3.pdf}}{%
  \phbox{#1}{#2}{#3}{#4}}}}}%
}

% Figura de evidencia completa:
% #1 ancho  #2 alto caja  #3 archivo  #4 descripcion  #5 caption
\newcommand{\evidence}[5]{%
  \begin{figure}[H]
  \centering
  \shot{#1}{#2}{#3}{#4}
  \caption{#5}
  \label{fig:#3}
  \end{figure}%
}

\captionsetup{font=small,labelfont=bf,skip=4pt}

% ---------------------------------------------------------------------

\title{\vspace{-1.4cm}\textbf{Workshop 6}\\[0.2em]
\large Testing \& CI/CD}
\author{Ismael Paredes \todo{other group members}\\[0.1em]
\normalsize Web Applications --- Semester II 2026 --- Prof. Francisco Hidrobo}
\date{\todo{submission date}}

\begin{document}

\maketitle
\vspace{-0.8cm}

% =====================================================================
\section{General Information}
% =====================================================================

\begin{table}[H]
\centering
\footnotesize
\begin{tabular}{@{}ll@{\hspace{1.5cm}}ll@{}}
\toprule
Group members & Ismael Paredes, \todo{others} & Python & \todo{version} \\
Course & Web Applications --- Sem. II 2026 & FastAPI & \todo{version} \\
Workshop & No. 6 --- Testing \& CI/CD & pytest & \todo{version} \\
Workshop date & October 30, 2026 & Ruff & \todo{version} \\
OS & Windows 11 Home (10.0.26200) & Hosting & Render (Free) \\
\bottomrule
\end{tabular}
\end{table}

\subsection*{Submission links}

\begin{table}[H]
\centering
\footnotesize
\begin{tabularx}{\textwidth}{@{}lX@{}}
\toprule
GitHub repository & \url{https://github.com/CipherTuring/Web-Applications} \\
Pull request & \todo{PR URL} \\
Actions run --- first passing run on the PR & \todo{run URL} \\
Actions run --- lint failure & \todo{run URL} \\
Actions run --- test failure & \todo{run URL} \\
Actions run --- passing run after fixes & \todo{run URL} \\
Actions run --- \texttt{main} after merge & \todo{run URL} \\
Public API URL & \todo{Render URL} \\
\bottomrule
\end{tabularx}
\end{table}

% =====================================================================
\section{Project and Technology Overview}
% =====================================================================

\todo{Short description of the API, the stack (FastAPI, Uvicorn, pytest +
TestClient, Ruff, GitHub Actions, Render) and why the API is self-contained.}

\begin{lstlisting}[style=plain]
TODO: workshop6/ directory tree
\end{lstlisting}

% =====================================================================
\section{Selected API Endpoint}
% =====================================================================

\todo{Which endpoint was selected (\texttt{GET /api/products}), what it
returns, and the small change made on the feature branch
(\texttt{GET /health}).}

\begin{lstlisting}[style=py]
# TODO: app/main.py
\end{lstlisting}

\begin{lstlisting}[style=py]
# TODO: tests/test_main.py
\end{lstlisting}

% =====================================================================
\section{CI Workflow}
% =====================================================================

\subsection{Workflow file}

\begin{lstlisting}[style=yml]
# TODO: .github/workflows/ci.yml
\end{lstlisting}

\subsection{Explanation of the checks}

\todo{Triggers (\texttt{pull\_request} and \texttt{push} to \texttt{main}),
the \texttt{lint} job (Ruff check + format check), the \texttt{test} job
(pytest), why they are independent jobs, and how the ruleset on
\texttt{main} turns them into required status checks.}

% =====================================================================
\section{Chronological Evidence}
% =====================================================================

% ---------------------------------------------------------------------
\subsection{Stage 1 --- Local change on a feature branch}
% ---------------------------------------------------------------------

\evidence{\textwidth}{5cm}{fig01-endpoint-code}{%
VS Code showing the change on branch \texttt{feature/health-endpoint}:
the new \texttt{GET /health} endpoint in \texttt{app/main.py} and its test
in \texttt{tests/test\_main.py} (diff view or both files side by side).}
{\todo{caption}}

\todo{Explanation of Figure~\ref{fig:fig01-endpoint-code}.}

\evidence{\textwidth}{5cm}{fig02-local-run}{%
Terminal running \texttt{uvicorn app.main:app --reload} with the request
log, next to the browser at \texttt{http://127.0.0.1:8000/health}
showing \texttt{\{"status": "ok"\}}.}
{\todo{caption}}

\todo{Explanation of Figure~\ref{fig:fig02-local-run}.}

\evidence{\textwidth}{5cm}{fig03-local-checks}{%
Terminal with the local checks passing: \texttt{ruff check .},
\texttt{ruff format --check .} and \texttt{pytest -v} with every test
\texttt{PASSED}.}
{\todo{caption}}

\todo{Explanation of Figure~\ref{fig:fig03-local-checks}.}

\evidence{\textwidth}{5cm}{fig04-git-push}{%
Terminal with \texttt{git checkout -b feature/health-endpoint},
\texttt{git add}, \texttt{git commit} and \texttt{git push -u origin
feature/health-endpoint}, including the push output.}
{\todo{caption}}

\todo{Explanation of Figure~\ref{fig:fig04-git-push}.}

% ---------------------------------------------------------------------
\subsection{Stage 2 --- Pull request and first CI run}
% ---------------------------------------------------------------------

\evidence{\textwidth}{5.5cm}{fig05-pr-opened}{%
GitHub pull request page \texttt{feature/health-endpoint} $\rightarrow$
\texttt{main}: title, description and the checks section with the
\texttt{lint} and \texttt{test} jobs.}
{\todo{caption}}

\todo{Explanation of Figure~\ref{fig:fig05-pr-opened}.}

\evidence{\textwidth}{5.5cm}{fig06-actions-pass}{%
GitHub Actions run summary for the PR: workflow name, triggering event
and commit, and both jobs (\texttt{lint}, \texttt{test}) in green, with
one job expanded showing its steps.}
{\todo{caption}}

\todo{Explanation of Figure~\ref{fig:fig06-actions-pass}.}

\evidence{\textwidth}{5.5cm}{fig07-branch-ruleset}{%
Repository Settings $\rightarrow$ Rules $\rightarrow$ Rulesets: ruleset
targeting \texttt{main} with ``Require a pull request'', ``Require status
checks to pass'' (\texttt{lint}, \texttt{test}) and ``Block force
pushes'' enabled.}
{\todo{caption}}

\todo{Explanation of Figure~\ref{fig:fig07-branch-ruleset}.}

% ---------------------------------------------------------------------
\subsection{Stage 3 --- Intentional failure 1: lint}
% ---------------------------------------------------------------------

\evidence{\textwidth}{5.5cm}{fig08-lint-fail-checks}{%
Pull request page after pushing the lint violation: \texttt{lint} job
failed (red), \texttt{test} job passed, and the merge box showing that
merging is blocked.}
{\todo{caption}}

\todo{Explanation of Figure~\ref{fig:fig08-lint-fail-checks}.}

\evidence{\textwidth}{5cm}{fig09-lint-fail-log}{%
Log of the failed \texttt{lint} job, expanded on the Ruff step: the
\texttt{F401} error with file, line and the unused import, and the
non-zero exit code.}
{\todo{caption}}

\todo{Explanation of Figure~\ref{fig:fig09-lint-fail-log}.}

\evidence{\textwidth}{4.5cm}{fig10-lint-fix-local}{%
Diff removing the unused import and the terminal with \texttt{ruff check .}
passing locally before pushing the fix.}
{\todo{caption}}

\todo{Explanation of Figure~\ref{fig:fig10-lint-fix-local}.}

% ---------------------------------------------------------------------
\subsection{Stage 4 --- Intentional failure 2: test}
% ---------------------------------------------------------------------

\evidence{\textwidth}{5.5cm}{fig11-test-fail-checks}{%
Pull request page after pushing the behaviour change: \texttt{lint}
passed, \texttt{test} failed, merging blocked.}
{\todo{caption}}

\todo{Explanation of Figure~\ref{fig:fig11-test-fail-checks}.}

\evidence{\textwidth}{5cm}{fig12-test-fail-log}{%
Log of the failed \texttt{test} job, expanded on the pytest step: the
failing test name and the \texttt{AssertionError} with expected vs.\ actual
value.}
{\todo{caption}}

\todo{Explanation of Figure~\ref{fig:fig12-test-fail-log}.}

\evidence{\textwidth}{5.5cm}{fig13-fix-pass-run}{%
GitHub Actions run for the fix commit: both jobs green again.}
{\todo{caption}}

\todo{Explanation of Figure~\ref{fig:fig13-fix-pass-run}.}

\evidence{\textwidth}{5.5cm}{fig14-pr-commit-history}{%
Commits tab (or timeline) of the pull request listing every commit in
order with its check status: pass, lint failure, lint fix, test failure,
test fix.}
{\todo{caption}}

\todo{Explanation of Figure~\ref{fig:fig14-pr-commit-history}.}

\subsection*{Failure summary}

\begin{table}[H]
\centering
\footnotesize
\begin{tabularx}{\textwidth}{@{}lXXX@{}}
\toprule
Failure & What was broken & What the log showed & How it was corrected \\
\midrule
Lint & \todo{} & \todo{} & \todo{} \\
Test & \todo{} & \todo{} & \todo{} \\
\bottomrule
\end{tabularx}
\end{table}

% ---------------------------------------------------------------------
\subsection{Stage 5 --- Merge to main}
% ---------------------------------------------------------------------

\evidence{\textwidth}{5.5cm}{fig15-pr-merged}{%
Pull request in the \emph{Merged} state, with all checks green and the
merge commit into \texttt{main}.}
{\todo{caption}}

\todo{Explanation of Figure~\ref{fig:fig15-pr-merged}.}

\evidence{\textwidth}{5.5cm}{fig16-main-run-pass}{%
GitHub Actions run triggered by the \texttt{push} event on \texttt{main}
(the merge commit), with both jobs green.}
{\todo{caption}}

\todo{Explanation of Figure~\ref{fig:fig16-main-run-pass}.}

% ---------------------------------------------------------------------
\subsection{Stage 6 --- Deployment}
% ---------------------------------------------------------------------

\evidence{\textwidth}{5.5cm}{fig17-render-config}{%
Render Web Service settings: repository and branch \texttt{main}, root
directory \texttt{workshop6}, build and start commands, Free instance and
Auto-Deploy set to ``After CI Checks Pass''.}
{\todo{caption}}

\todo{Explanation of Figure~\ref{fig:fig17-render-config}.}

\evidence{\textwidth}{5.5cm}{fig18-render-deploy-live}{%
Render Events / deploy log: deploy triggered by the \texttt{main} merge
commit, build and start output, and the service status \emph{Live}.}
{\todo{caption}}

\todo{Explanation of Figure~\ref{fig:fig18-render-deploy-live}.}

\evidence{\textwidth}{5cm}{fig19-public-endpoint}{%
Browser at the public URL \texttt{https://<service>.onrender.com/api/products}
(and/or \texttt{/health}) returning the JSON response.}
{\todo{caption}}

\todo{Explanation of Figure~\ref{fig:fig19-public-endpoint}.}

% =====================================================================
\section{CI vs.\ Continuous Delivery vs.\ Continuous Deployment}
% =====================================================================

\todo{Definition of each practice and how this pipeline maps to them:
Actions (CI), always-deployable \texttt{main} (Delivery), Render
auto-deploy after checks pass (Deployment), and the ruleset as the gate.}

% =====================================================================
\section{Conclusions and Limitations}
% =====================================================================

\subsection*{Conclusions}
\todo{Conclusions.}

\subsection*{Limitations}
\todo{Limitations.}

\end{document}
